% ============================================================
% II. NGHIÊN CỨU LIÊN QUAN
% ============================================================
\section{Các nghiên cứu liên quan}

\begin{frame}{Nghiên cứu liên quan: dữ liệu và baseline}
  \begin{columns}
    \column{0.49\textwidth}
      \textbf{II-A. FoodSeg103 \& fine-grained segmentation}
      \begin{itemize}
        \item 7.118 ảnh, 103 lớp, 42.097 mask \cite{benchmarkfood}
        \item DeepLabv3+ chỉ đạt \textbf{34.2 mIoU} trên benchmark này \cite{benchmarkfood,chen2017deeplab}
        \item ReLeM giúp SeTR tăng 41.3 $\to$ 43.9 mIoU \cite{benchmarkfood,zheng2021setr}
        \item Yêu cầu: nhạy với texture, boundary, class imbalance
      \end{itemize}
      \column{0.50\textwidth}
      \textbf{II-B. Lightweight segmentation}
      \begin{itemize}
        \item BiSeNet: 68.4\% mIoU @ 105 FPS \cite{bisenet}
        \item Fast-SCNN: 68.0\% @ 123.5 FPS, 1.11M params \cite{poudel2019fastscnn}
        \item PIDNet-S: 78.6\% @ 93.2 FPS \cite{xu2023pidnet}
      \end{itemize}
    

  \end{columns}
    \begin{block}{Nhận xét nhanh}
        \begin{itemize}
          \item Bài toán food segmentation khó hơn benchmark Cityscapes.
          \item Mô hình nhẹ đạt tốc độ tốt, nhưng thường yếu ở chi tiết biên và lớp hiếm.
        \end{itemize}
      \end{block}
\end{frame}

\begin{frame}{Nghiên cứu liên quan: cơ chế tăng cường chi tiết}
  \begin{columns}
    \column{0.52\textwidth}
      \textbf{II-C. Cơ chế tăng cường local detail}
      \begin{itemize}
        \item Gated-SCNN: shape stream riêng \cite{takikawa2019gatedscnn}
        \item PIDNet: 3 nhánh detail--context--boundary \cite{xu2023pidnet}
        \item PointRend: tinh chỉnh tại điểm bất định \cite{kirillov2020pointrend}
        \item Deep-TEN: texture như tín hiệu thống kê \cite{zhang2017deepten}
      \end{itemize}

        Các hướng trên đều thêm tín hiệu \textbf{hình học cục bộ} hoặc \textbf{texture} vào pipeline segmentation.
    \column{0.44\textwidth}
    \textbf{II-D. Khoảng trống nghiên cứu}\\
        Chưa có pipeline baseline đủ sạch để kiểm tra hệ thống các cơ chế detail-aware trên FoodSeg103
  \end{columns}
        \begin{block}{Câu hỏi nghiên cứu của nhóm}
        Nếu chuẩn hóa baseline trước, thì những lỗi nào thực sự do kiến trúc lightweight gây ra và nên sửa ở đâu?
      \end{block}
\end{frame}
